% ============================================================================
% style.tex — the identity of this reconstructed book (Spivak, Calculus).
%
% Author-specific macros live here and GROW after every chapter. Packages
% and figure style live in the kit; do not add a capability that the next
% book will need into this file.
%
% Design decisions (current):
%   * Font: Computer Modern (the classic Spivak identity — chosen by the
%     owner over matching the 2008 Times re-typeset). No font package is
%     loaded; CM is the LaTeX default. Figures use the document font.
%   * Class: book, 11pt, A4, oneside, openany, moderate margins (~2.2cm).
%     The original is a print-optimised book with huge margins that look
%     tiny on a tablet; we target a layout that reads well on BOTH screen
%     and paper. oneside/openany: no forced blank rectos.
%   * Structure: \frontmatter (title–contents, roman pages) until Chapter 1,
%     \mainmatter (Parts I–V, arabic from 1), \backmatter (reading, answers,
%     glossary, index).
%   * Theorem voice: Spivak labels environments in full capitals —
%     THEOREM / DEFINITION / PROOF — statement in italics, proof upright.
%     Numbered per chapter, no chapter prefix ("THEOREM 1", not "THEOREM 1.1").
%   * Properties: the twelve "basic properties" are labelled (P1)..(P12).
%   * Problems: numbered 1., 2., ...; starred (*) for hard ones; sub-parts
%     use small roman (i)(ii) or letters (a)(b) depending on the problem.
% ============================================================================

% --- headers/footers -------------------------------------------------------
\pagestyle{fancy}
\fancyhf{}
\fancyhead[L]{\nouppercase{\leftmark}}
\fancyhead[R]{\nouppercase{\rightmark}}
\fancyfoot[C]{\thepage}
\renewcommand{\headrulewidth}{0.4pt}
\setlength{\headheight}{14pt}

% Book class owns the chapter counter. Theorems/equations reset per chapter
% and display their plain number (matching the book's "THEOREM 1", "(1)", ...).
\numberwithin{equation}{chapter}
\renewcommand{\theequation}{\arabic{equation}}

% ============================================================================
% Chapter opener. The classic Spivak chapter begins on a fresh page with a
% large "CHAPTER n" and the title beneath it. #1 is accepted for call-site
% compatibility and ignored — the book counter is the number.
% ============================================================================
\newcommand{\spivakchapter}[2]{%
  \clearpage
  \refstepcounter{chapter}%
  \chaptermark{#2}%
  \thispagestyle{plain}%
  \vspace*{3cm}%
  \begin{center}%
    {\scshape\Large Chapter}\\[0.6em]%
    {\fontsize{60}{70}\selectfont\bfseries\thechapter}\\[0.6em]%
    {\Large\bfseries #2}%
  \end{center}%
  \vspace*{3cm}%
  \clearpage
}

% Part divider (e.g. "PART I / PROLOGUE"). The source uses a sparse page
% with the part name; we reproduce that sparseness instead of filling it.
\newcommand{\partdivider}[2]{%
  \clearpage
  \thispagestyle{empty}%
  \vspace*{6cm}%
  \begin{center}%
    {\scshape\LARGE Part #1}\\[1.2em]%
    {\Huge\bfseries #2}%
  \end{center}%
  \vspace*{\fill}%
  \clearpage
}

% Epigraph page (a short quote opening a part/chapter). Kept short and
% centred, exactly as in the source — do not reflow following text into it.
\newcommand{\epigraph}[2]{%
  \clearpage
  \thispagestyle{empty}%
  \vspace*{\fill}%
  \begin{flushright}%
    \begin{minipage}{0.6\textwidth}%
      \itshape #1\par
      \vspace{1em}%
      {\upshape\scshape #2}%
    \end{minipage}%
  \end{flushright}%
  \vspace*{\fill}%
  \clearpage
}

% ============================================================================
% Theorem-style environments in Spivak's voice: bold capitals label, no
% trailing period (the book prints "THEOREM 1", not "THEOREM 1."), statement
% in italics, proof body upright and closed with a solid box.
% ============================================================================
\newtheoremstyle{spivak}%        name
  {12pt}%                        space above
  {12pt}%                        space below
  {\itshape}%                    body font
  {}%                            indent
  {\bfseries}%                   head font
  {}%                            head punct (none)
  {5pt}%                         space after head
  {}%                            head spec

\newtheoremstyle{spivakdef}%     name (upright body, for definitions)
  {12pt}%
  {12pt}%
  {\normalfont}%
  {}%
  {\bfseries}%
  {}%
  {5pt}%
  {}%

\theoremstyle{spivak}
\newtheorem{theorem}{THEOREM}[chapter]
\newtheorem{lemma}{LEMMA}[chapter]
\newtheorem{corollary}{COROLLARY}[chapter]
\renewcommand{\thetheorem}{\arabic{theorem}}
\renewcommand{\thelemma}{\arabic{lemma}}
\renewcommand{\thecorollary}{\arabic{corollary}}

\theoremstyle{spivakdef}
\newtheorem{definition}{DEFINITION}[chapter]
\renewcommand{\thedefinition}{\arabic{definition}}

% Proof: Spivak prints "PROOF" in capitals (no period) and closes with a box.
\renewcommand{\proofname}{PROOF}
\renewcommand{\qedsymbol}{\ensuremath{\blacksquare}}
\makeatletter
\renewenvironment{proof}[1][\proofname]{%
  \par\pushQED{\qed}\normalfont\topsep6\p@\@plus6\p@\relax
  \trivlist
  \item[\hskip\labelsep\bfseries #1]\ignorespaces
}{%
  \popQED\endtrivlist\@endpefalse
}
\makeatother

% ============================================================================
% Properties (P1)..(P12). \prop{N} prints the label; the statement follows.
% ============================================================================
\newcommand{\prop}[1]{\textbf{(P#1)}}

% ============================================================================
% Problems. \pnum{<label>} typesets a problem number in bold at the left
% margin; the label is written verbatim so starred problems read "*15.".
% ============================================================================
\newcommand{\pnum}[1]{%
  \par\medskip\noindent
  \hangindent=2.4em\hangafter=1%
  \makebox[2.4em][l]{\textbf{#1}}%
  \ignorespaces
}

% Sub-part lists: small roman "(i)(ii)" and letters "(a)(b)".
\newlist{romans}{enumerate}{1}
\setlist[romans]{label={(\roman*)},leftmargin=2.6em,itemsep=0pt,topsep=2pt}
\newlist{letters}{enumerate}{1}
\setlist[letters]{label={(\alph*)},leftmargin=2.6em,itemsep=0pt,topsep=2pt}
